% APA 7 style preamble for pandoc -> LaTeX -> PDF.
% preamble.lua splices this into header-includes ahead of the paper's own.

% URLs: break only after a slash, so "https://" never splits after the colon.
\usepackage{xurl}
\makeatletter\renewcommand*{\UrlBigBreaks}{\do\/}\makeatother

% Page number top right, no rule. APA 7 puts it in the header; a paper that needs
% it in the footer sets `pagenumber: bottomright` in its front matter.
\usepackage{fancyhdr}
\pagestyle{fancy}
\fancyhf{}
\ifdefined\rbPageNumberBottom
  \fancyfoot[R]{\thepage}
\else
  \fancyhead[R]{\thepage}
\fi
\renewcommand{\headrulewidth}{0pt}

% The professional title page also carries a running head: the short title in
% capitals, flush left. preamble.lua fills \rbShortTitle from the paper's
% `shorttitle:` field, or from its title when there is none.
\providecommand{\rbShortTitle}{}
\ifdefined\rbProfessionalTitlePage
  \fancyhead[L]{\rbShortTitle}
\fi

% Double and single spacing. Pandoc's template loads setspace only when the
% paper sets `linestretch`, but the title page and the table bodies need it
% either way, so the style asks for it itself. A second load is a no-op.
\usepackage{setspace}

% Keep a block together when there is room, else start it on the next page.
\usepackage{needspace}

% Reserve room for a table so it is not split across a page break. #1 is the
% number of lines the table itself occupies (rows plus its rules), #2 the
% number of double-spaced body lines above it (its label and title, or 0 when
% it has none). The table is set \small and single-spaced, so its line height
% is measured inside that group rather than guessed from the body's.
%
% A table taller than a page still splits, and longtable repeats its header
% row, which is what APA and MLA both ask for.
\newlength{\rbBodyLine}
\newcommand{\rbTableNeed}[2]{%
  \setlength{\rbBodyLine}{\baselineskip}%
  \begingroup
    \small\singlespacing
    \needspace{\dimexpr #2\rbBodyLine + #1\baselineskip\relax}%
  \endgroup}

% APA 7 title page. Everything 12pt and double-spaced: title in bold three lines
% down, then one blank line, then the author block. The title page is page 1 and
% carries the page number.
%
% The student format (the default) closes with the due date, and its author
% block runs author, affiliation, course and instructor. The professional
% format, selected with `titlepage: professional`, drops the due date and adds
% the running head above.
\usepackage{etoolbox}
\makeatletter
\ifdefined\rbProfessionalTitlePage
  \renewcommand{\maketitle}{%
    \thispagestyle{fancy}%
    \begin{center}%
      \setstretch{2}\selectfont
      \vspace*{3\baselineskip}
      {\normalsize\bfseries\@title\par}
      \vspace{\baselineskip}
      {\normalsize\@author\par}
    \end{center}%
    \newpage}
\else
  \renewcommand{\maketitle}{%
    \thispagestyle{fancy}%
    \begin{center}%
      \setstretch{2}\selectfont
      \vspace*{3\baselineskip}
      {\normalsize\bfseries\@title\par}
      \vspace{\baselineskip}
      {\normalsize\@author\par}
      {\normalsize\@date\par}
    \end{center}%
    \newpage}
\fi
\makeatother

% A paper that supplies its own title page sets `titlepage: false`.
\ifdefined\rbNoTitlePage\renewcommand{\maketitle}{}\fi
